\section{Hauptverstärker}

Der Hauptverstärker dient dazu, das Signal des Vorverstärkers (PMT) auf eine Praktisch messbare Amplitude zu verstärken und die Signalform so zu formen, dass sie für die weitere Verarbeitung praktisch ist. Hierbei muss die Pulsdauer deutlich reduziert werden um die Überlappung mehrerer Pulse zu seltener zu machen, und die Pulse werden auf eine (ungefähre) Gaussform gebracht um das Signal zu Rauschen Verhältniss zu verbessern\cite[280]{leotechniques}.

Zur Pulsformung wird üblicherweise ein so genanntes \textit{CR-RC Netzwerk} genutzt. Dabei wird das Signal durch einen Tiefpass und einen Hochpassfilter gefiltert \cite[281]{leotechniques}.

Es wird ein \texttt{ORTEC 572A Amplifier} genutzt, welcher diese Funktionen erfüllt \cite{shaping_amp}.
